\chapter*{Conclusion générale et perspectives}
\addcontentsline{toc}{chapter}{Conclusion générale et perspectives}
\markboth{Conclusion générale et perspectives}{}

Ce projet avait pour but de connaître le vent partout et à toutes les hauteurs utiles à partir de
mesures de surface éparses. Nous avons conçu, entraîné, évalué et mis en service PC-FNO3D, un
opérateur neuronal de Fourier tridimensionnel contraint par la physique.

\paragraph{Bilan.}
Sur le jeu synthétique, le modèle de référence reconstruit en 91\,ms un champ de 1,5 million de
nœuds, respecte la conservation de la masse ($2{,}6\cdot10^{-5}$), reste robuste avec un quart
des capteurs et divise par 7,5 à 9,8 l'erreur maximale des méthodes d'interpolation classiques.
Transposé au Grand Tunis, il apprend une correction du profil ERA5 : le modèle livré atteint une
erreur globale de 0,38\ms, garantit la conservation de la masse par construction, fournit une
incertitude calibrée (88\,\% de couverture pour l'intervalle à 80\,\%), fait mieux qu'ERA5 sur les
radiosondages et le vaut aux stations au sol. Il est servi par une API sécurisée, un serveur MCP
et l'application trilingue ATOM Wind, avec le temps réel et la prévision à trois jours.

\paragraph{Limites.}
Le niveau le plus bas (25\,m) reste à 0,53\ms, au-dessus du seuil ; le champ réel est un peu trop
lisse (rugosité 0,805) ; aux stations, la qualité est bornée par la vérité générée. L'étude OS6 a
montré que le modèle livré ne tire aucun gain de capteurs ajoutés, ce qui empêche de dimensionner
le réseau avec lui. Ces limites sont mesurées et documentées.

\paragraph{Apports personnels.}
\acompleter{Compétences acquises : apprentissage profond scientifique, physique de l'atmosphère,
ingénierie des données météorologiques, MLOps sur GPU distant, développement web et API, conduite
d'expériences avec règles de décision pré-enregistrées.}

\paragraph{Perspectives.}
\begin{itemize}
  \item \textbf{un modèle sensible aux capteurs} : réentraîner avec des capteurs simulés denses et
        tirés au hasard à partir de la vérité de surface, pour que le modèle exploite les mesures
        ANEO, puis refaire l'étude OS6 ;
  \item \textbf{la prévision OS7} : terminer l'affinage sur GFS et le juger contre GFS brut et la
        persistance ;
  \item \textbf{les données ANEO réelles} : brancher les premiers mâts Porteur et recaler le plan
        du réseau sur leurs mesures ;
  \item \textbf{l'état atmosphérique complet} : reconstruire aussi la température, l'humidité et la
        pression, déjà présentes en entrée ;
  \item \textbf{l'hébergement} : un service permanent au lieu d'une session GPU à la demande, et le
        test utilisateur prévu (UX1, VIS1) avec cinq personnes non spécialistes.
\end{itemize}
